\chapter{Louis Dumont on Caste: Purity, Impurity and Hierarchy}

\section{Dumont as Indologist and Structuralist}

\begin{KeyConcept}
\begin{itemize}
\item \textbf{Louis Dumont} is a French \textbf{Indologist / Orientalist} --- he relied on \textbf{ancient Sanskritic texts} to understand caste (he even learnt Sanskrit).
\item \textbf{Indology is about WHAT you read} (ancient Sanskritic texts), not HOW you read them. Dumont's distinctive contribution is the method he used: he read the Indological texts through \textbf{structuralism} (the theoretical perspective of \textbf{Levi-Strauss}).
\item \textbf{Structuralism} reads the world through \textbf{binary opposites} (good/bad, rich/poor, hot/cold, life/death): you cannot make sense of one category without its opposite. These oppositional categories are \textbf{cognitive categories} --- they structure how we make sense of everyday life.
\item Dumont therefore looked for the \textbf{binary opposites within the ancient Sanskritic texts} on caste.
\end{itemize}
\end{KeyConcept}

\section{Bouglé's Three Features and the Underlying Principle}

\begin{KeyConcept}
\begin{itemize}
\item Before Dumont, the French scholar \textbf{Celestine Bouglé} (Dumont's intellectual predecessor, himself influenced by Levi-Strauss) had identified \textbf{three features of caste}:
\item 1. \textbf{Hierarchy} --- ritual hierarchy (status-based; the top is ritually pure, the bottom ritually impure).
\item 2. \textbf{Hereditary occupation and specialization} --- trade secrets passed through caste (caste itself acted like a 'patent').
\item 3. \textbf{Repulsion (separation)} --- castes maintain distance from each other, through \textbf{endogamy} and \textbf{commensality} (restrictions on whom to take food/water from), e.g. the Agraharam/Kudiyana/Cheri separation.
\end{itemize}
\end{KeyConcept}

\begin{KeyConcept}
\begin{itemize}
\item \textbf{Dumont's key insight}: beneath these three features (the superstructures) lies one \textbf{underlying base} --- the \textbf{binary opposition of purity and impurity}.
\item Hierarchy has the pure at top and impure at bottom; occupation is divided into \textbf{pure occupations} (priestly ritual) and \textbf{impure occupations} (manual scavenging, menial tasks); repulsion separates the pure castes from the impure ones.
\item So the \textbf{underlying principle governing the features of caste is the binary opposition of purity and impurity}.
\end{itemize}
\end{KeyConcept}

\begin{KeyConcept}
This is the sociologist's task --- \textbf{debunking / revealing what is hidden} (Peter Berger's idea of sociology as debunking; Robert Merton's 'latent' functions): every sociologist reveals something hidden, and Dumont revealed the base beneath the superstructure.
\end{KeyConcept}

\section{Purity and Impurity as the Ideology of Caste}

\begin{KeyConcept}
\begin{itemize}
\item Dumont believes \textbf{caste is the central feature of Hindu society}, established on the principle of purity and impurity.
\item But \textbf{purity and impurity go beyond caste} --- they govern everything in Hindu life: no haircut on Tuesday, Navratri fasts (no chicken), \textbf{silk is pure while cotton is impure}, death brings 12 days of impurity, childbirth brings impurity, menstruation is ritually impure (the \textbf{Sabarimala} case), etc.
\item Even in non-vegetarian food there is purity/impurity: for a Bengali, \textbf{fish is pure} (offered to the goddess) but chicken is impure.
\item The consciousness of the Hindu is shaped by purity/impurity --- it is about \textbf{inequality} (pure/impure), not equality.
\end{itemize}
\end{KeyConcept}

\begin{KeyConcept}
\begin{itemize}
\item Dumont's propositions (to write down):
\item \textbf{The ideology of the caste system is based on the binary opposition of purity and impurity.}
\item \textbf{Caste is a set of economic, political and social relationships sustained by religious values.}
\item \textbf{Social supremacy in Indian society is closely associated with ritual status, articulated as the oppositional unity of pure and impure.}
\item \textbf{The rules of pure and impure are central to the establishment and functioning of the caste system.}
\end{itemize}
\end{KeyConcept}

\section{The Texts: P. V. Kane and Harita}

\begin{KeyConcept}
\begin{itemize}
\item To decipher the purity/impurity principle, Dumont referred to two key texts: \textbf{P. V. Kane's 'History of Dharmashastra'} and \textbf{Harita}.
\item \textbf{Harita} spoke of \textbf{three kinds of purity}: (1) \textbf{Kula} (clan/family --- descent from an ancestor), (2) \textbf{Sharira} (the body), and (3) \textbf{objects of everyday use}.
\item \textbf{P. V. Kane} held that the purity of a caste is \textbf{maintained through endogamy}.
\end{itemize}
\end{KeyConcept}

\section{Status over Power: The Supremacy of the Religious}

\begin{KeyConcept}
\begin{itemize}
\item In India, \textbf{ritual status is more important than secular power}: the \textbf{king employs priests; priests do not employ kings}.
\item Dumont's line: \textbf{'In India, status is more important than power'} --- it does not matter whether you are an IAS officer or the Prime Minister; what matters to the Hindu is \textbf{who is ritually pure (status-wise)}.
\item The notion of purity/impurity comes from the \textbf{supremacy given to religious values in Indian society}; ritual supremacy is accorded more importance than secular power.
\item (\textbf{T. N. Madan} and \textbf{Ashis Nandy} similarly say religion is 'in our blood' --- so how can one talk of secularism?)
\end{itemize}
\end{KeyConcept}

\begin{KeyConcept}
If everyone --- even kings and low castes --- follows the purity/impurity rituals made by Brahmins, then everyone has been \textbf{Sanskritized}, and the way one thinks is \textbf{Brahminical}; hence superiority goes to the Brahmins.
\end{KeyConcept}

\section{India vs the West: Homo Hierarchicus vs Homo Aequalis}

\begin{KeyConcept}
\begin{itemize}
\item Dumont compares Indian and Western society in his classic book \textbf{'Homo Hierarchicus'} (the West being 'European/American'), again using \textbf{binary opposites}:
\item \textbf{Hierarchy vs Equality} --- Indian society is \textbf{hierarchical} (based on purity/impurity); Western society is \textbf{egalitarian} (no purity/impurity).
\item \textbf{Holism vs Individualism} --- in India \textbf{man does not exist} (only caste/society exists; caste rules regulate you); in the West the \textbf{individual exists} (free, governed only by free speech).
\item In India, the individual exists only as the \textbf{Sanyasi (renouncer)} --- who is \textbf{casteless} ('don't ask a sadhu his jati').
\item Dumont calls Hindu society \textbf{'Homo Hierarchicus'} and the West \textbf{'Homo Aequalis'}.
\end{itemize}
\end{KeyConcept}

\section{Homo Major and Homo Minor: Holism vs Individualism}

\begin{KeyConcept}
\begin{itemize}
\item Dumont coined two further terms: \textbf{'Homo Major'} (India) and \textbf{'Homo Minor'} (the West).
\item \textbf{Homo Major} --- in India, \textbf{society prevails over the individual} (man is a \textbf{collective being}; e.g. an unhappy couple does not divorce because 'what will society say?'). Indians are characterised more by \textbf{holism, religion and hierarchy} --- religious principles are more important than the political/economic domain.
\item \textbf{Homo Minor} --- in the West, \textbf{the individual prevails over society}. The values governing Western people are \textbf{equality/egalitarianism, free market economy, a highly rational political system, democracy, individualism and secularism} (complete separation of state and religion --- you are a \textbf{citizen first, Christian second}; in India you are Hindu first).
\item The \textbf{exception} proves each rule: in India, individualism appears only through the \textbf{renouncer (Sanyasi)}; in the West, society-over-individual appears only in \textbf{totalitarianism / hyper-nationalism} (e.g. Hitler's German consciousness).
\end{itemize}
\end{KeyConcept}

\begin{KeyConcept}
\begin{itemize}
\item Dumont is describing the \textbf{traditional Hindu society} (the Indological version, via texts including the Dharma Shastras) and the \textbf{post-French-Revolution Western society} --- he is not saying one is good and the other bad, only that their \textbf{governing values differ}.
\item Dumont greatly admired \textbf{Ambedkar} --- who also held that \textbf{democracy and Hinduism are opposite to each other}.
\end{itemize}
\end{KeyConcept}

\section{Temporary vs Permanent Impurity}

\begin{KeyConcept}
\begin{itemize}
\item Dumont distinguished \textbf{two types of impurity}:
\item \textbf{Temporary impurity} --- the victim is the \textbf{upper caste} (no haircut on certain days, death and birth impurities). It can be \textbf{purified through rituals} --- for the upper caste, pollution is \textbf{like anomie} (a temporary phase, fixed by purification).
\item \textbf{Permanent impurity} --- the victim is the \textbf{lower caste}: their occupations (menial tasks) are permanently polluted, and their impurity \textbf{cannot be purified by any ritual}.
\item Hence untouchability is permanent for the lower castes, while the upper castes can always 'purify' themselves.
\end{itemize}
\end{KeyConcept}

\section{Dumont's Field Work: The Pramalai Kallar}

\begin{KeyConcept}
\begin{itemize}
\item Dumont is \textbf{both Indological and empirical} (integrating both, unlike Ghurye who was largely Indological): he also did \textbf{field work}.
\item His field-study book is on the \textbf{'Pramalai Kallar'} sub-caste (Madurai, Tamil Nadu) --- \textbf{'A South Indian Sub-Caste: Social Organization and Religion of the Pramalai Kallar'}.
\item In \textbf{Homo Hierarchicus} he used the \textbf{Varna model}; in the field study he used the \textbf{Jati model}.
\item Even in the field he found the \textbf{binary opposition of purity/impurity} being followed --- residential separation, ritual separation, endogamy and food restrictions among the Pramalai Kallar --- proving that Hindu religious values govern the cognition of all Hindus, North or South.
\end{itemize}
\end{KeyConcept}

\section{The Binary Oppositions of the Caste Structure}

\begin{KeyConcept}
\begin{itemize}
\item Dumont pushed binary opposition to its extreme in describing the caste structure, as a chain of oppositions:
\item \textbf{Savarna vs Avarna} (Brahmin/Kshatriya/Vaishya/Shudra vs the untouchables).
\item \textbf{Caste vs non-caste} (the four varnas vs the untouchables / Ati-Shudra / Adi Dravida).
\item \textbf{Dvija vs Ekaj} (twice-born B,K,V vs the once-born Shudra).
\item \textbf{Priestly vs non-priestly} (Brahmin vs K,V,S).
\item \textbf{Warrior vs non-warrior} (Kshatriya vs B,V,S).
\item \textbf{Dominant vs non-dominant} (Brahmin vs K,V,S --- even the dominant castes, after gaining secular power, still claim Kshatriya status).
\item \textbf{Brahmin vs non-Brahmin} (B vs K,V,S --- the same opposition Beteille found).
\end{itemize}
\end{KeyConcept}

\begin{KeyConcept}
\textbf{Conclusion}: Dumont believed Indian society is based on \textbf{cognitive structures} which have their origin in \textbf{religious consciousness, rooted in history} (the Indological texts). His approach is therefore called the \textbf{cognitive-historical approach}.
\end{KeyConcept}

\section{Substantialization of Caste}

\begin{KeyConcept}
\begin{itemize}
\item On \textbf{social change in caste} (the modern avatar of caste --- as Ghurye, Srinivas and Beteille also discussed), Dumont argues:
\item Earlier castes were \textbf{interdependent} (specialization and division of labour); today they are \textbf{no longer interdependent} --- every caste has become a \textbf{self-sufficient, impenetrable block}, identical to and in competition with one another (caste associations, vote banks, e.g. the All-India Thakur Association).
\item He calls this process the \textbf{'substantialization of caste'} --- the same phenomenon Ghurye called \textbf{'caste patriotism'} and Srinivas called the \textbf{'vote bank'}.
\item In traditional India there was \textbf{integration between castes} (so he becomes a Durkheimian); in modern India there is \textbf{conflict between castes} (so he becomes a Marxist).
\end{itemize}
\end{KeyConcept}

\begin{QuickRevision}
\begin{itemize}
\item Dumont = Indologist (Sanskrit texts) + structuralist (Levi-Strauss); reads Indological texts through binary opposites.
\item Bouglé's three features (hierarchy, hereditary occupation, repulsion) $\rightarrow$ Dumont's underlying base = binary opposition of purity/impurity.
\item Purity/impurity governs all Hindu life (haircut, Navratri, silk vs cotton, death/birth, menstruation/Sabarimala); caste = central feature of Hindu society.
\item Texts: P. V. Kane (History of Dharmashastra; purity via endogamy) and Harita (kula/sharira/everyday objects).
\item Ritual status > secular power ('kings employ priests'); 'status is more important than power'; religious supremacy.
\item India vs West: hierarchy vs equality; holism vs individualism; Homo Hierarchicus vs Homo Aequalis; Homo Major (society over individual) vs Homo Minor (individual over society); renouncer = only Indian individual; totalitarianism = only Western holism.
\item Temporary impurity (upper caste, purifiable, like anomie) vs permanent impurity (lower caste, never purified).
\item Field work: Pramalai Kallar (Madurai) --- Jati model; purity/impurity found even in the field.
\item Binary oppositions: Savarna/Avarna; caste/non-caste; Dvija/Ekaj; priestly/non-priestly; warrior/non-warrior; dominant/non-dominant; Brahmin/non-Brahmin.
\item Conclusion: cognitive structures from religious consciousness rooted in history = cognitive-historical approach.
\item Substantialization of caste: interdependence $\rightarrow$ self-sufficient competing blocks (like Ghurye's caste patriotism, Srinivas's vote bank).
\end{itemize}
\end{QuickRevision}

